% ID: SEL-D0077-OA-E01
% Tanda S077.
\textbf{a)} Despejamos $y$ en las tres restricciones:
\[
 y\le4+\frac{x}{2},\qquad y\ge13-x,\qquad y\le49-4x.
\]
En el primer cuadrante, el recinto queda por encima de la recta decreciente $y=13-x$ y por debajo de las otras dos. Es el triángulo cerrado representado, incluidos sus lados.
\begin{center}
\setlength{\unitlength}{1mm}
\begin{picture}(110,76)(-10,-9)
\put(0,0){\vector(1,0){98}}\put(0,0){\vector(0,1){64}}
\put(99,0){$x$}\put(1,64){$y$}
{\color{gray}
\qbezier(0,20)(45,38.75)(90,57.5)
\qbezier(6,60)(42,30)(78,0)
\qbezier(55.5,60)(64.5,30)(73.5,0)}
\linethickness{0.45mm}
\qbezier(36,35)(48,40)(60,45)
\qbezier(60,45)(66,25)(72,5)
\qbezier(72,5)(54,20)(36,35)
\put(23,36){\scriptsize $(6,7)$}
\put(59,48){\scriptsize $(10,9)$}
\put(73,5){\scriptsize $(12,1)$}
\put(48,28){\scriptsize Recinto}
\end{picture}
\end{center}

\textbf{b)} Intersecamos las fronteras por parejas. De $2y-x=8$ y $x+y=13$ resulta $3y=21$, luego $(6,7)$. De $2y-x=8$ y $y+4x=49$ resulta $(10,9)$. De $x+y=13$ y $y+4x=49$ resulta $(12,1)$. Los tres puntos satisfacen las restricciones restantes. Ningún punto de los ejes es factible, por lo que estos son todos los vértices.

\textbf{c)} Una función lineal alcanza sus extremos en algún vértice de un polígono cerrado y acotado. Evaluamos:
\[
 F(6,7)=2,\qquad F(10,9)=6,\qquad F(12,1)=44.
\]
El mínimo es $\boxed{2}$ y se alcanza únicamente en $(6,7)$. El máximo es $\boxed{44}$ y se alcanza únicamente en $(12,1)$.
